% Owner: Abhinav (P8.5); numbers from docs/study_area.md and the feature table (C4)
\section{Study area}\label{sec:study-area}

The study area is \textbf{Bengaluru South}, a \SI{582}{\km\squared} rectangle from
\ang{77.52}\,E to \ang{77.74}\,E and \ang{12.72}\,N to \ang{12.94}\,N
(Figure~\ref{fig:study-area}). It covers the south-eastern growth corridor along Hosur Road
(Electronic City, Bommasandra), the Kanakapura Road fringe in the south-west, and
\textbf{Bannerghatta National Park}, of which \SI{41.3}{\km\squared} lies inside the area. All
analysis is done in UTM zone 43N (EPSG:32643) on a grid of 58{,}218 cells of
\SI{100}{\m}.

\paragraph{Why this area.} Fourteen candidate areas of 400--\SI{620}{\km\squared} around Indian
cities were compared on their land-cover mix, their persistent growth since 2018 and their
OpenStreetMap coverage (\texttt{docs/study\_area.md}). Bengaluru South was chosen for four
reasons:
\begin{itemize}
  \item \textbf{strong, persistent growth:} 8.4\,\% of the land that was non-built in 2018 was
        built in both 2022 and 2023, which gives enough change to validate against;
  \item \textbf{a clear forest block} (Bannerghatta) next to cropland and scrub, which gives the
        three-class map a real forest class;
  \item \textbf{mostly gentle terrain:} elevations of \SIrange{727}{1022}{\m} (median
        \SI{898}{\m}) and a median cell slope of \ang{2.5}; under 1\,\% of the cells are
        steeper than \ang{15}, mostly in the Bannerghatta hills;
  \item \textbf{good OpenStreetMap coverage in 2018:} about 151{,}000 building footprints
        were already mapped on 1 January 2018. Several other candidates (Coimbatore,
        Bhubaneswar) were mapped in bulk only later, so they could not supply buildings as of
        the baseline date, which the temporal validation needs.
\end{itemize}

\paragraph{Land cover.} In 2018 the ESRI map classes 53.6\,\% of the area as built, 18.4\,\% as
crops, 18.1\,\% as rangeland, 7.2\,\% as trees and 2.5\,\% as water. By 2023 built land had
risen to 61.5\,\%, mostly at the expense of cropland (12.2\,\%). The built-up core lies in the
north, toward central Bengaluru; growth is concentrated in the south and east. Many small lakes
(tanks) are scattered across the area, and some of them are overgrown with vegetation, which the
exclusion mask has to handle.

\begin{figure}[ht]
  \centering
  \includegraphics[width=0.78\textwidth]{map_lulc_2023}
  \caption{Study area: ESRI land cover in 2023 \citep{karra2021}, with the area outline.
           Bannerghatta National Park is the tree and rangeland block in the south-west.}
  \label{fig:study-area}
\end{figure}
